\documentclass[10pt,a4paper]{amsart}
\usepackage[margin=19mm]{geometry}
\usepackage{amsmath,amssymb,mathtools}
\usepackage[scr=rsfs,cal=euler]{mathalfa}
\usepackage{xfrac}
\usepackage[unicode,hidelinks,bookmarks=false]{hyperref}
\newcommand{\ie}{i.e.\ }
\newcommand{\cf}{cf.\ }
\newcommand{\resp}{respectively\ }
\newcommand{\Subsection}{\subsection}
\providecommand{\nobreakdash}{\nobreak\hbox{-}}
\setlength{\emergencystretch}{2em}
\multlinegap0pt
\usepackage{color}
\newtheorem{Theorem}{Theorem}
\newtheorem{lemma}{Lemma}
\newtheorem{corollary}{Corollary}
\newtheorem{proposition}{Proposition}
\newtheorem{remark}{Remark}
\newtheorem{Definition}{Definition}
\newtheorem*{ack}{Acknowledgement}
 \newtheorem{ex}{Example}
 \numberwithin{equation}{section}
\usepackage{multirow}
\usepackage{mathrsfs}
\usepackage{appendix}
\usepackage{textcomp}
\usepackage{manyfoot}
\usepackage{booktabs}
\usepackage{algorithm}
\usepackage{algorithmicx}
\usepackage{algpseudocode}
\usepackage{listings}
\usepackage{upref}
\usepackage{float}
\newcommand{\RR}{{\mathbb R}} \newcommand{\DD}{{\mathbb D}}
\newcommand{\BB}{{\mathbb B}} \newcommand{\CC}{{\mathbb C}}
\newcommand{\ZZ}{{\mathbb Z}} \newcommand{\QQ}{{\mathbb Q}}
\newcommand{\NN}{{\mathbb N}}
\def\a{\nu} \def\vp{\mathcal{T}} \def\msk{\medskip} \def\ol{\overline}
\def\bege{\begin{equation}} \def\ende{\end{equation}} \def\bsk{\bigskip}
\def\cent{\centerline} \def\a{\nu} \def\om{\omega}
\def\b{\beta} \def\g{\mathcal{T}}\def\ol{\overline} \def\ve{\varepsilon}
\def\d{\delta} \def\pt{\partial} \def\bin{\atopwithdelims} \def\begr{\begin{eqnarray}}
\def\endr{\end{eqnarray}} \def\qand{\quad\mbox{ and }\quad}
\def\qfor{\quad\mbox{ for }\quad} \def\qie{\quad\mbox{ i.e. }\quad}
\def\bege{\begin{equation}} \def\ende{\end{equation}}
\def\begr{\begin{eqnarray}} \def\endr{\end{eqnarray}}
\def\bnum{\begin{enumerate}} \def\enum{\end{enumerate}}
\allowdisplaybreaks[4]
\begin{document}
\title[Commutants of complex symmetric weighted composition operato]{Commutants of complex symmetric weighted composition operators on Fock space}
\author{Aakriti Sharma}
\date{}
\maketitle
\noindent\emph{Source and licence.} Commutants of complex symmetric weighted composition operators on Fock space. Version arXiv:2606.25544v1 (CC BY 4.0). This material is distributed under the Creative Commons Attribution 4.0 International licence, http://creativecommons.org/licenses/by/4.0/, as recorded in the arXiv OAI licence metadata for this exact version and shown on the abstract page https://arxiv.org/abs/2606.25544v1. Full rights notice: licenses/arxiv-2606.25544-CC-BY-4.0.txt.\par\smallskip
\noindent\textit{Editorial scope.} Counted unit: the original \texttt{Theorem} environment \texttt{t3} at source lines 289--297 together with its complete proof at lines 298--367; the delivered document keeps source lines 191--368 so that every referenced label and every citation resolves inside it. Native editable source is the paper's own concatenated TeX; the AMS wrapper is editorial formatting only. No publisher class, upstream script or unrelated artwork is redistributed.
\begin{Theorem}\cite{PVL}\label{t1}
  Let $W_{\psi, \varphi}$ be bounded weighted composition operator on $\mathcal{F}^{2}$ induced by two entire function $\psi$ and $\varphi$ with $\psi\not\equiv 0 $. Then  $W_{\psi, \varphi}$ is $J_{a,b,c}$-symmetric  if and only if
  $$\varphi(z)= Az+B,\quad \psi(z)=Ce^{Dz}$$ 
  with $C\neq 0$, $D=aB-bA+b$ where $A$, $B$ are complex constants satisfying either $|A|< 1$
  or $|A|=1, D+A\overline{B}=0$.
\end{Theorem}
To know more about the boundedness, compactness, self-adjointness and normality of weighted composition operators on Fock space, we refer \cite{TL1}and references therein.
The following eigenstructure result will also be needed.
\begin{Theorem}\cite{PVL}\label{t2}Let $W_{\psi, \varphi}$ be a $J_{a,b,c}$-symmetric bounded weighted composition operator on $\mathcal{F}^{2}$ induced by two entire function $\psi$ and $\varphi$ with $\psi\not\equiv 0 $. If $\varphi$ has a fixed point $d(=\frac{B}{1-A})$, then 
$$\Omega_{k}(z)= (z-d)^{k}K_{\overline{\alpha}}(z),\;\;\; k\in \mathbb{N}\cup \{0\}, \alpha=\frac{D}{1-A}.  $$
are eigenvectors of $W_{\psi, \varphi}$ corresponding to the eigenvalues $A^{k}\psi(d)$.
\end{Theorem}


\section{Commutants of complex symmetric weighted composition operators}\label{sec:prelimi}
In this section, we investigate the commutant of a complex symmetric weighted composition operator. We derive necessary and sufficient conditions for two weighted composition operators to commute. The commuting relation leads to functional equations involving the symbols, and these equations are analyzed using reproducing kernel techniques and analytic function theory. To establish the main theorem, we first require the following auxiliary lemmas.
\begin{lemma}\label{1l}
    Let $W_{f,g}$ and  $W_{\psi, \varphi}$ be two weighted composition operators on $\mathcal{F}^{2}$, where $f, g, \psi, \varphi$ are entire functions. Then $W_{f,g}\in \{W_{\psi, \varphi}\}^{'}$ if and only if $f.\psi \circ g =\psi .f\circ \phi$ and $\varphi \circ g = g \circ \varphi$.
\end{lemma}
\begin{proof}
    We compute
    $$W_{f,g}W_{\psi, \varphi}= W_{f.\psi \circ g , \varphi \circ g} ~\text{ and }~ W_{\psi, \varphi}W_{f,g}= W_{\psi. f\circ \varphi , g\circ \varphi}. $$ Since we have
    $$W_{f,g}\in \{W_{\psi, \varphi}\}^{'} \text{ if and only if } W_{f,g}W_{\psi, \varphi}= W_{\psi, \varphi}W_{f,g}.$$
    Hence $W_{f,g}\in \{W_{\psi, \varphi}\}^{'}$
    if and only if $f.\psi \circ g =\psi .f\circ \phi$ and $\varphi \circ g = g \circ \varphi$.
\end{proof}
\begin{lemma}\label{l2} 
Let $f$ and $g$ be two entire function with $f\neq 0$ and $W_{\psi, \varphi}$ be $J_{a,b,c}$-complex symmetric weighted  composition operators on $\mathcal{F}^{2}$ with $\varphi$ has a fixed point 
$0\neq d(=\frac{B}{1-A})$ in $\mathbb{C}$. If $W_{f,g} \in \{W_{\psi, \varphi}\}^{'}$, then  $\varphi$ and $g$ have a same fixed point $d$.
\end{lemma}
\begin{proof}
   Since by lemma\ref{1l}, we have 
   $$W_{f,g} \in \{W_{\psi, \varphi}\}^{'} ~\text{ if and only if }~ f.\psi \circ g =\psi .f\circ \phi ~\text{ and }~ \varphi \circ g = g \circ \varphi $$
   %$W_{f,g}W_{\psi, \varphi}=W_{\psi, \varphi}W_{f,g}$ 


   and by Theorem \ref{t1} $W_{\psi, \varphi}$ is $J_{a,b,c}$-complex symmetric weighted  composition operators with $\varphi$ has a fixed point 
$0\neq d(=\frac{B}{1-A})$ in $\mathbb{C}$. Then we have 
$$\varphi(g(d)) = g (\varphi(d))= g(d)$$
which implies $g(d)=d$. This proves that $\varphi$ and $g$ have a same fixed point $d$.
\end{proof}

\begin{lemma}\label{l1}Let $\eta= \alpha(1-\beta)$, $\alpha=\frac{D}{1-A}$, $d=\frac{B}{1-A}$ and $\beta \in \mathbb{C}$, Then the following  identities hold:
\begin{enumerate}
\item $D+\alpha A = \alpha$ and $\alpha B= dD$
\item $\eta+D\beta= D+\eta A$ 
\item $D(1-\beta)d= \eta B$ 
\item $ A(1-\beta)d+B= (1-\beta)d+B\beta$
\end{enumerate}
\end{lemma}
\begin{proof}
    
       \begin{enumerate}
           \item $D+\alpha A= D+\frac{D}{1-A}A=\frac{D}{1-A}=\alpha$ and $\alpha B= \frac{D}{1-A} B= dD$\\
       \item We compute
       \begin{align*}
            \eta+D\beta &= \alpha(1-\beta)+D\beta\\
            &=\frac{D}{1-A}(1-\beta)+D\beta\\
            &= \frac{D-DA\beta}{1-A}
        \end{align*}  
        %\end{minipage}
        and
       % \begin{minipage}{.5\linewidth}
            \begin{align*}
            D+\eta A &= D+\alpha (1-\beta)A\\
            &= D+\frac{D}{1-A}(1-\beta)A\\
            &=\frac{D-DA\beta}{1-A}
        \end{align*}
        %\end{minipage}\\
        
        Therefore $\eta+D\beta= D+\eta A$ holds.\\
        \item  We compute
        \begin{align*}
        D(1-\beta)d &= \frac{D(1-\beta)B}{1-A}\\
        &= \alpha(1-\beta)B\\
        &= \eta B 
        \end{align*}
    Therefore $D(1-\beta)d= \eta B$ holds.\\
    
    \item  We compute
   % \begin{minipage}{.3\linewidth}
        \begin{align*}
        A(1-\beta)d+B &= \frac{A(1-\beta)B}{1-A}+B\\
        &= \frac{B(1-AB)}{1-A}
    \end{align*}
    %\end{minipage}
    and 
    %\begin{minipage}{.3\linewidth}
       \begin{align*}
        (1-\beta)d+B\beta &= (1-\beta)\frac{B}{1-A}+ B\beta\\
        &= \frac{B(1-AB)}{1-A}
    \end{align*} 
   % \end{minipage}\\
    
Therefore $ A(1-\beta)d+B= (1-\beta)d+B\beta$ holds.
\end{enumerate}
\end{proof}


\begin{Theorem}\label{t3}
Let $f$ and $g$ be two entire function with $f\neq 0$ and $W_{\psi, \varphi}$ be $J_{a,b,c}$-complex symmetric weighted  composition operators on $\mathcal{F}^{2}$ with $\varphi$ has a fixed point 
$d(=\frac{B}{1-A})$ in $\mathbb{C}$. Then $W_{f,g} \in \{W_{\psi, \varphi}\}^{'}$ if and only if symbols functions $f$ and $g$ are of the following forms:
\begin{equation}\label{e1}
  g(z)= \beta z+ (1-\beta) d \text{~~ and ~~} f(z)= f(d) e^{\eta(z-d)}  
\end{equation}

where  $\beta \in \mathbb{C}$ and $ \eta = \alpha (\beta -1)$ .
\end{Theorem}

\begin{proof}
First, by using identity $(1)$ of lemma\ref{l1}, we compute
\begin{align*}
    W_{\psi, \varphi}K_{\overline{\alpha}}(z)&= C.e^{Dz}e^{\alpha \varphi(z)}\\
    &= C. e^{(D+\alpha A)z+\alpha B}\\
    &= \psi(d)K_{\overline{\alpha}}(z)
\end{align*}
    

  Since $W_{f,g} \in \{W_{\psi, \varphi}\}^{'}$ and $W_{\psi, \varphi}$ be $J_{a,b,c}$-complex symmetric , then by theorem\ref{t2} we have

  \begin{align*}
  W_{\psi, \varphi}W_{f,g}K_{\overline{\alpha}}&= W_{f,g}W_{\psi, \varphi}K_{\overline{\alpha}}\\
  &= \psi(d)W_{f,g}K_{\overline{\alpha}}
  \end{align*}
If $W_{f,g}K_{\overline{\alpha}}=0$, then $K_{\overline{\alpha}}$ is an eigen vector of $W_{f,g}$ corresponding to eigen value 0.
If $W_{f,g}K_{\overline{\alpha}}\neq 0$, then $W_{f,g}K_{\overline{\alpha}}$ is an eigen vector of $W_{f,g}$ corresponding to eigen value 
$\psi(d)$. Then by theorem\ref{t1} we have $W_{f,g}K_{\overline{\alpha}}= \gamma_{1}K_{\overline{\alpha}}$ for some non zero $\gamma_{1}\in \mathbb{C}$. Again by theorem\ref{t1}, the function $\Omega_{1}= (z-d)K_{\overline{\alpha}}$ is an eigen vector of $W_{\psi, \varphi}$ corresponding to eigen value $A\psi(d)$.
Since $W_{f,g} \in \{W_{\psi, \varphi}\}^{'}$, we have 
\begin{align*}W_{\psi, \varphi}W_{f,g}\Omega_{1}&=W_{f,g}W_{\psi, \varphi}\Omega_{1}\\
&=A\psi(d) W_{f,g}\Omega_{1} 
\end{align*}

This show that $W_{f,g}\Omega_{1}$  is an eigenvector corresponding to the eigenvalue $A\psi(d)$. Therefore $W_{f,g}\Omega_{1}=\gamma_{2}\Omega_{1}$, for some $\gamma_{2} \in \mathbb{C}$. Thus we have

\begin{align*}
   W_{f,g}\Omega_{1}&=\gamma_{2}\Omega_{1}\\
    f(z) (g(z)-d)K_{\overline{\alpha}}(g(z))&= \gamma_{2} (z-d)K_{\overline{\alpha}}(z)\\
    (W_{f,g}K _{\overline{\alpha}}(z))(g(z)-d))&= \gamma_{2}(z-d)K_{\overline{\alpha}}(z)\\
    \gamma_{1}K_{\overline{\alpha}}(z)(g(z)-d)&=\gamma_{2} (z-d)K_{\overline{\alpha}}(z)\\
    g(z)&=\beta(z-d)+d\\
    g(z)&=\beta z+ (1-\beta) d
\end{align*}
where $\beta=\frac{\gamma_{2}}{\gamma_{1}} \in \mathbb{C}.$

Now, let $z=d$ in  $W_{f,g}K_{\overline{\alpha}}(z)= \gamma_{1}K_{\overline{\alpha}}(z)$, by lemma\ref{l2}, we get that $\gamma_{1}=f(d)$. Therefore, we have
\begin{align*}
    W_{f,g}K_{\overline{\alpha}}(z)&= \gamma_{1}K_{\overline{\alpha}}(z)\\
   f(z)e^{\alpha g(z)}&=f(d)e^{\alpha z} \\
   f(z)&= f(d)e^{\alpha z-\alpha g(z)}\\
   f(z)&= f(d) e^{\alpha z-\alpha(\beta z+ (1-\beta) d)}\\
   f(z)&= f(d) e^{\eta(z-d)}
\end{align*}
where $\eta = \alpha(1-\beta)$.\\
Conversely, suppose that $f$ and $g$ are of the form in \eqref{e1}. To complete the proof we have to show that 
%$W_{f,g}W_{\psi, \varphi}=W_{\psi, \varphi}W_{f,g}$ implies
$$f.\psi \circ g = \psi. f\circ \varphi \text{ and }\varphi \circ g= g\circ \varphi.$$
Thus we have 
\begin{align*}
 f(z)\psi (g(z)) &= f(d) e^{\eta(z-d)} C. e^{Dg(z)}\\
 &= C.f(d)e^{\eta z-\eta d+D(\beta z +(1-\beta)d)}\\
 &= C.f(d)e^{(\eta+D\beta)z+(D(1-\beta)d-\eta d)}
\end{align*}
and 
\begin{align*}
    \psi(z) f( \varphi (z))&= C. e^{Dz} f(d)e^{\eta [(Az+B)-d]}\\
    &= C. e^{Dz} f(d) e^{(D+\eta A)z+(\eta B-\eta d)}
\end{align*}
Clearly, by using identities $(1)$ and $(2)$ of lemma\ref{l1}, we conclude that $f.\psi \circ g = \psi. f\circ \varphi $. Now we compute
\begin{align*}
    \varphi( g(z))&= Ag(z)+B\\
    &= A(\beta z+ (1-\beta) d)+B
\end{align*}
and 
\begin{align*}
    g(\varphi(z))&= \beta \varphi(z)+ (1-\beta) d\\
    &= \beta (Az+B)+(1-\beta) d\\
\end{align*}
Again by using identity $(3)$ of lemma\ref{l1}, we conclude that $\varphi \circ g= g\circ \varphi.$ This completes the proof of the theorem.
\end{proof}


\begin{thebibliography}{99}
\bibitem[PVL]{PVL} { P.~V.~Hai and L.~H.~Khoi: }{\it Complex symmetry of weighted composition operators on the Fock space, } { Journal of Mathematical Analysis and Applications,} {433 (2)},{1757--1771}, (2016) {https://doi.org/10.1016/j.jmaa.2015.08.069}.
\bibitem[TL1]{TL1} { T. Le: }{\it Normal and isometric weighted composition operators on the fock space,} { Bull. Lond. Math. Soc.,} (46(4),{847--856}, (2014)
\end{thebibliography}
\end{document}
